% ===========================================================================
% STANDALONE TEST FILE -- compile this on its own to preview the figure.
%
%   pdflatex test_fig_pooling.tex
%
% The page geometry below gives a ~12 cm text width, matching the Springer
% Nature single-column layout, so \resizebox{\textwidth} scales the figure by
% the same factor it will in the real manuscript. A wider test page would make
% the figure look larger than it will on submission.
%
% To move this into the manuscript, copy ONLY the \begin{figure} ... \end{figure}
% block below, and add these two lines to the manuscript preamble:
%
%   \usepackage{tikz}
%   \usetikzlibrary{calc,arrows.meta,decorations.pathreplacing}
% ===========================================================================

\documentclass[11pt]{article}

% page geometry chosen to mimic the journal's text width
\usepackage[a4paper,left=4.5cm,right=4.5cm,top=2.5cm,bottom=2.5cm]{geometry}

% maths -- needed for \mathbf, \sum, \sigma, \Rightarrow in the node labels
\usepackage{amsmath}
\usepackage{amssymb}

% the figure itself
\usepackage{graphicx}      % provides \resizebox
\usepackage{xcolor}        % provides \textcolor (also pulled in by tikz)
\usepackage{tikz}
\usetikzlibrary{calc,arrows.meta,decorations.pathreplacing}

\begin{document}

%% NOTE: the caption cross-references Fig.~\ref{fig:architecture}, which does not
%% exist in this standalone file, so it will typeset as "Fig. ??". That is expected
%% here and resolves once the figure is placed in the manuscript.

\begin{figure}[htbp]
\centering
\resizebox{\textwidth}{!}{%
\begin{tikzpicture}[
  font=\sffamily\footnotesize,
  main/.style  = {draw, rounded corners=2pt, minimum width=4.6cm,
                  minimum height=1.0cm, align=center, inner sep=2pt},
  feat/.style  = {main, fill=orange!14, draw=orange!75!black},
  les/.style   = {main, fill=blue!12,   draw=blue!60},
  att/.style   = {main, fill=blue!24,   draw=blue!70},
  vec/.style   = {main, fill=green!12,  draw=green!45!black},
  outp/.style  = {main, fill=violet!14, draw=violet!65!black},
  stop/.style  = {draw=red!65!black, line width=1pt, rounded corners=2pt,
                  fill=red!12, minimum width=4.6cm, minimum height=0.9cm,
                  align=center, inner sep=2pt},
  side/.style  = {draw, rounded corners=2pt, minimum width=2.0cm,
                  minimum height=1.0cm, align=center, inner sep=2pt,
                  fill=green!10, draw=green!45!black},
  loss/.style  = {draw, rounded corners=2pt, minimum width=2.4cm,
                  minimum height=1.0cm, align=center, inner sep=2pt,
                  fill=black!5, draw=black!55},
  note/.style  = {font=\sffamily\scriptsize, text=black!65, align=left},
  rnote/.style = {font=\sffamily\scriptsize, text=red!55!black, align=left},
  hdr/.style   = {font=\sffamily\small\bfseries, text=black!80},
  ar/.style    = {-{Stealth[length=2.4mm]}, semithick, draw=black!65},
  dash/.style  = {-{Stealth[length=2.4mm]}, semithick, draw=black!55, dashed}
]

%% ================= main column ============================================
\node[feat] (m0) at (5.6,   0.00) {Encoder feature map $\mathbf{F}$\\[-1pt]
                                   {\scriptsize $512 \times 16 \times 16$}};
\node[les]  (m1) at (5.6,  -1.50) {Lesion decoder};
\node[les]  (m2) at (5.6,  -2.90) {Lesion logits $\mathbf{L}$\\[-1pt]
                                   {\scriptsize $512 \times 512$}};
\node[les]  (m3) at (5.6,  -4.20) {$\sigma(\mathbf{L})$ --- predicted lesion\\[-1pt]
                                   {\scriptsize probability map}};
\node[stop] (m4) at (5.6,  -5.40) {\textbf{\textcolor{red!65!black}{STOP-GRADIENT}}\\[-1pt]
                                   {\scriptsize detach from the computation graph}};
%% NOTE: keep each math expression on ONE source line. A $...$ group broken
%% across lines inside a TikZ node can be typeset literally instead of as maths.
\node[att]  (m5) at (5.6,  -6.80) {Resample to $16\times16$, normalise\\[-1pt]
                                   {\scriptsize to unit sum $\Rightarrow$ attention $\mathbf{A}$}};
\node[vec]  (m6) at (5.6,  -8.20) {$\mathbf{v}_{\mathrm{lesion}} = \sum_{h,w} \mathbf{A}_{h,w}\,\mathbf{F}_{:,h,w}$\\[-1pt]
                                   {\scriptsize $512$-d}};
\node[vec]  (m7) at (5.6,  -9.80) {$[\,\mathbf{v}_{\mathrm{global}} \,\|\, \mathbf{v}_{\mathrm{lesion}}\,]$\\[-1pt]
                                   {\scriptsize concatenation, $1024$-d}};
\node[outp] (m8) at (5.6, -11.20) {Local-pathology head\\[-1pt]
                                   {\scriptsize linear $1024 \rightarrow 9$ logits}};

\draw[ar] (m0) -- (m1);
\draw[ar] (m1) -- (m2);
\draw[ar] (m2) -- (m3);
\draw[ar] (m3) -- (m4);
\draw[ar] (m4) -- (m5);
\draw[ar] (m5) -- (m6);
\draw[ar] (m6) -- (m7);
\draw[ar] (m7) -- (m8);

%% ================= global-pooling branch (left) ===========================
\node[side] (l1) at (0.8, -6.80) {Global average\\pooling};
\node[side] (l2) at (0.8, -8.20) {$\mathbf{v}_{\text{global}}$\\[-1pt]
                                  {\scriptsize $512$-d}};

\draw[ar] (3.30, 0.25) -- (0.80, 0.25) -- (0.80, -6.30);
\draw[ar] (l1) -- (l2);
\draw[ar] (0.80, -8.70) -- (0.80, -9.80) -- (3.30, -9.80);

%% ================= feature re-use channel =================================
\draw[dash] (3.30, -0.25) -- (2.40, -0.25) -- (2.40, -8.20) -- (3.30, -8.20);
\node[note, rotate=90, anchor=south] at (2.22, -4.20)
      {same $\mathbf{F}$, weighted by $\mathbf{A}$};

%% ================= segmentation loss (right) ==============================
\node[loss] (dice) at (9.50, -2.90) {Dice loss\\[-1pt]
                                     {\scriptsize sole supervision}};
\draw[ar] (7.90, -2.90) -- (dice.west);

%% ================= stop-gradient annotation ===============================
\node[rnote, anchor=west, text width=2.6cm] at (8.30, -5.40)
      {No gradient returns from the classification loss to the lesion decoder.};

\node[note, anchor=west, text width=2.6cm] at (8.30, -11.20)
      {Multi-label output: nine independent sigmoids.};

%% ================= headers ================================================
\node[hdr] at (5.6, 1.05) {Lesion-guided pooling};
\node[hdr, text=green!35!black] at (0.8, 1.05) {global branch};

\end{tikzpicture}}
\caption{Lesion-guided pooling. The multi-label pathology classifier is given localised evidence
by reusing the lesion decoder's own output as a spatial attention map. The predicted lesion
probability map is \emph{detached} from the computation graph, resampled to the encoder's spatial
resolution and normalised to unit sum, yielding attention weights $\mathbf{A}$. These weights
pool the encoder features into a lesion-focused descriptor
$\mathbf{v}_{\text{lesion}}$, which is concatenated with the conventional globally averaged
descriptor $\mathbf{v}_{\text{global}}$ before classification. The stop-gradient is essential:
without it the classification loss could reshape the lesion decoder to produce convenient
pseudo-lesions, improving the classifier at the expense of the segmentation output a clinician
sees. With it, the lesion decoder remains supervised solely by its own Dice term and the
attention is read-only.}\label{fig:pooling}
\end{figure}
\end{document}
